\documentclass[10pt,a4paper]{amsart}
\usepackage[margin=19mm]{geometry}
\usepackage{amsmath,amssymb,mathtools}
\usepackage[scr=rsfs,cal=euler]{mathalfa}
\usepackage{xfrac}
\usepackage[unicode,hidelinks,bookmarks=false]{hyperref}
\newcommand{\ie}{i.e.\ }
\newcommand{\cf}{cf.\ }
\newcommand{\resp}{respectively\ }
\newcommand{\Subsection}{\subsection}
\providecommand{\nobreakdash}{\nobreak\hbox{-}}
\setlength{\emergencystretch}{2em}
\multlinegap0pt
\usepackage{tikz}
\usetikzlibrary{cd}
\usepackage{amssymb}
\usepackage[shortlabels]{enumitem}
\usepackage{xcolor}
\usepackage{tikz}
\newtheorem{proposition}{Proposition}
\newcommand{\GL}{\mathrm{GL}}
\newcommand{\SL}{\mathrm{SL}}
\newcommand{\ZZ}{\mathbb{Z}}
\newcommand{\ZZhat}{\widehat{\ZZ}}
\title{Supersingular isogeny graphs and Hecke modules with level structure}
\author{Leonardo Col\`o, David Kohel}
\date{Source: arXiv:2604.00184v1 (CC BY 4.0)}
\begin{document}
\maketitle

\noindent\emph{Source and licence.} Supersingular isogeny graphs and Hecke modules with level structure. Version arXiv:2604.00184v1 (CC BY 4.0), distributed by arXiv under the Creative Commons Attribution 4.0 International licence, http://creativecommons.org/licenses/by/4.0/, as recorded in the arXiv licence metadata for this exact version and shown on the abstract page https://arxiv.org/abs/2604.00184v1. Full licence notice: \texttt{licenses/arxiv-2604.00184-CC-BY-4.0.txt}.\par\smallskip

\noindent\textit{Editorial scope.} Counted unit: the article's proposition proposition (source0.tex lines 812--820). The article's complete original proof of it is delivered with it (source0.tex lines 822--842, one \texttt{proof} environment of 21 lines). No mathematics is written, repaired or re-derived: the statement and the proof are the article's own lines, in the article's own notation, and no retained line is substituted or re-flowed.  No further source-local context is needed: every cross-reference of every retained line resolves inside the two delivered spans.  Nothing was omitted from any retained span, so the package carries no omission record.  No retained line contains a citation, so the wrapper carries no bibliography and no bibliography entry.  The retained lines contain the article's own diagram code, which the wrapper typesets with the standard tikz packages; no image file is delivered and the package declares no asset dependency.  Editorial formatting only, in addition: an AMS \texttt{amsart} preamble that restates the article's own declarations and macros used by the retained lines, namely the environments \texttt{proposition} on separate counters, together with the standard packages those lines need.  The retained environments print in this document as Proposition 1 in order of appearance, whereas the article prints the same items with the numbering of its own full document.  The complete native source of the article as distributed in the arXiv e-print for this exact version is bundled unchanged as \texttt{sources/197571/source0.tex}.
\begin{proposition}
  Let $H_1$ and $H_2$ be open supgroups of $G = \GL_2(\ZZhat)$. Any two of following
  conditions implies the third.
  \begin{enumerate} %itemize}
    \item The subgroups $H_1$ and $H_2$ are independent in $G$.
    \item The subgroups $H_1$ and $H_2$ are geometrically independent in $G$.
    \item The subgroups $\det(H_1)$ and $\det(H_2)$ are independent in $\ZZhat^*$.
  \end{enumerate} %itemize}
\end{proposition}

\begin{proof}
  For any open subgroup $G \subset \GL_2(\ZZhat)$ of level $N$,
  and reduction map $\pi_N: \GL_2(\ZZhat) \to \GL_2(\ZZ/N\ZZ)$,
  we set $G_0 = \pi_N(G)$ and $G_1 = G_0 \cap \SL_2(\ZZ/N\ZZ)$.
  From the reduction $\bmod N$ of the exact sequence,
  $$
  \begin{tikzcd}
    1 \arrow[r] &
    \SL_2(\ZZhat) \cap G \arrow[r] \arrow[d] &
    G \arrow[r] \arrow[d] &
    \det(G) \arrow[r] \arrow[d] & 1 \\
    1 \arrow[r] & G_1 \arrow[r] & G_0 \arrow[r] & \det(G_0) \arrow[r] & 1
  \end{tikzcd}
  $$
  and the identity $[\GL_2(\ZZhat):G] = [\GL_2(\ZZ/N\ZZ):G_0]$, the
  multiplicative relation
  $$
  [\GL_2(\ZZ/N\ZZ):G_0] = [\SL_2(\ZZ/N\ZZ):G_1][(\ZZ/N\ZZ)^*:\det(G_0)].
  $$
  gives the required dependency relation between the three independence conditions.
\end{proof}

\end{document}
